 \subsection{Algorithm for \textsc{Weighted-Frac-Match}()}

The final ingredient we need for \textsc{WeightedM-or-E*}(), is the following lemma, which computes a fractional matching on low capacity edges.

\begin{lemma2}
Consider a weighted bipartite graph $G$ with edge capacity function $\kappa$, and edge weights in $\set{1,2,\cdots,W}$. There is an algorithm that in $O(\nicefrac{m}{\eps}\cdot \log \nicefrac{1}{\eps})$ time finds a fractional matching $\vec{x}$ obeying the capacity function $\kappa$ such that $\sum_{e\in E}w(e)\cdot x(e)\geqslant (1-\eps)\cdot \textsf{mwm}(G,\kappa)$. 
\end{lemma2}

In order to come up with this algorithm, first recall the maximum weight capacitated fractional matching linear program, and its corresponding dual.

\begin{equation}
\begin{array}{ll@{}ll}
\text{maximize}  & \displaystyle\sum\limits_{e\in E} w(e)\cdot x(e)  \\
\text{subject to}& \displaystyle x(e)\leqslant \kappa(e)& & \forall e\in E \\
& \displaystyle \sum_{e\ni v} x(e) \leqslant 1 & & \forall v\in V
\end{array}
\end{equation}

\begin{equation}
\begin{array}{ll@{}ll}
\text{minimize}  & \displaystyle\sum\limits_{v\in L} y(v) +\sum\limits_{u\in R} y(u) +\sum_{e\in E} z(e)\cdot \kappa(e) &  \\
\text{subject to}& \displaystyle yz(e)\geqslant w(e) &  \forall e\in E 
\end{array}
\end{equation}

Here, $yz(e)=y(u)+y(v)+z(e)$, where $u,v$ are the two endpoints of $e$. We will come up with an algorithm that maintains the following invariants at all times. We denote $\delta_0=\eps\cdot W$, and $\delta_{i}=\nicefrac{\delta_0}{2^i}$. Additionally, for an edge $e$, we denote by $w_i(e)=\left\lfloor \frac{w(e)}{\delta_i}\right \rfloor\cdot \delta_i$, the truncated weight function at scale $i$. We additionally assume that $W$ and $\eps$ are powers of $2$. 

\begin{definition2}
Given a bipartite graph $G$, with a capacity function $\kappa$ on the edges of the graph. Suppose $\vec{x}$ is a fractional matching on edges of the graph, obeying the capacity function $\kappa$. We define $G_x$ to be the residual graph with respect to $\vec{s}$. This is a directed graph, where, corresponding to every edge $e$, there is a forward edge $e_f$, and a backward edge $e_b$, and capacity of $e_f$ is $\kappa(e)-x(e)$, and $e_b$ has capacity $x(e)$. 
\end{definition2}

 We additionally define free vertices as follows:

\begin{definition2}
Given a graph $G$ and a fractional matching $\vec{x}$, we say that a vertex $v$ is free with respect to $\vec{x}$ if $\sum_{e\ni v}x(e)<1$.
\end{definition2}

\begin{property}\label{prop:invariant} The algorithm will maintain at all times, a fractional matching $\vec{x}$, and dual variables $y(u)$ for all $u\in V$, and $z_e$ for all $e\in E$. The algorithm will proceed in $i=1,2,\cdots, \log W$ scales, and in each scale $i$, we will maintain the following properties in the residual graph $G_x$. 
\begin{enumerate}
\item \textbf{Granularity:} The dual variable $y(u)$ is a multiple of $\delta_i$, and similarly, $z(e)$ is a multiple of $\delta_i$. \label{item:gran}
\item \textbf{Near Domination:} For all edges $e_f$ and $e_b$, we have $yz(e_f)\geqslant w_i(e)-\delta_i$, and $yz(e_b)\geqslant w_i(e)-\delta_i$. \label{item:domin}
\item \textbf{Tightness:} For all back edges $e_b$, we have, $yz(e_b)\leqslant w_i(e)+2\cdot(\delta_j-\delta_i)$, where $j$ is the scale at which $e_b$ last existed in $G_x$ (with non-zero capacity), and was never removed. We say that in this case that $e_b$ is a backward edge of type $j$.  \label{item:tightness}
\item \textbf{Free Duals:} The duals $y(u)$ for free $u\in L$ are equal, and are at most the dual values for all other vertices. We stop scale $i$ when free vertex duals in $L$ reach value $\nicefrac{W}{2^{i+1}}-\nicefrac{\eps\cdot W}{2^{i+1}}$. On the other hand, the duals $y(u)$ for all free $u\in R$ are zero throughout the run of the algorihm. At the end of the scale $\log W$, all free vertex duals $y(u)$ are $0$.\label{item:freeduals}
\item \textbf{Complementary Slackness:} If $z(e)>0$ for some $e\in E$, then $x(e)=\kappa(e)$. \label{item:slackness}
\end{enumerate} 
\end{property}

Our goal will be to first show that if an algorithm maintains invariants mentioned in \Cref{prop:invariant}, then at the end of the algorithm, we have a matching $M$ with $w(M)\geqslant (1-\eps)\cdot \textsf{mwm}(G,\kappa)$. In order to do this, we start with some definitions.

\begin{definition2}[Eligible Edges]
Let $\vec{x}$ be the current fractional matching maintained by the graph, and let $G_x$ denote the residual graph with respect to $\vec{x}$. In scale $i$, we define a forward edge $e_f$ to be eligible if $yz(e_f)=w(e_f)-\delta_i$, and backward edge $e_b$ is eligible if $yz(e_b)-w(e_b)$ is a non-negative multiple of $\delta_i$. We use $G^t_{x}$ to denote the eligible subgraph of $G_x$. 
\end{definition2}
We first show the following intermediate claim, and thereafter, we show this claim combined with \Cref{prop:invariant}, shows that at the of the algorithm, we have $\sum_{e\in E} w(e)\cdot x(e)\geqslant (1-\eps)\cdot \textsf{mwm}(G,\kappa)$. 

\begin{claim2}\label{claim:forward}
All eligible forward edges in scales $0$ to $i$ have weights at least $\frac{W}{2^{i+1}}$. Moreover, in scale $i$, for all backward edges of type $j$, we have $yz(e)\leqslant (1+4\eps)\cdot w(e)$.
\end{claim2}
\begin{proof}
First, consider forward edges $e_f$, these are eligible if they have the property that $yz(e)=w(e)-\delta_i$ in scale $i$. Note that the $yz(e)$ values are lower bounded by the $y(u)$ value of the free vertices which is at least $\nicefrac{W}{2^i}-\delta_i$. Therefore, it follows that $w(e)\geqslant \nicefrac{W}{2^i}$. This shows the first part of the claim. To show the second part of the claim, since $e$ is a backward edge of type $j$, this means that before scale $j$, it was solely a forward edge. Consequently, we have that $w(e)\geqslant \nicefrac{W}{2^{j+1}}$. To prove the second part of the argument, we have the following line of reasoning due \Cref{prop:invariant}\ref{item:tightness} and discussion above, 
\begin{align*}
yz(e)-w(e)&\leqslant 2(\delta_j-\delta_i)\\
&< 2\cdot \delta_j\\
&\leqslant 8\cdot\eps\cdot w(e)
\end{align*}
This proves the claim.
\end{proof}

Now, we want to a bound on the weight of the fractional matching output by the algorithm. 

\begin{claim2}
Let $\vec{x}$ be the fractional matching returned by the graph. Then, we have, $\sum_{e\in E}x(e)\cdot w(e)\geqslant (1-\eps)\cdot \textsf{mwm}(G)$.
\end{claim2}
\begin{proof}
In order to see this, consider the dual variables returned by the algorithm. We know that,
\begin{align*}
\sum_{u\in V}y(u)+\sum_{e\in E}z(e)\cdot \kappa(e)\geqslant (1-\eps)\cdot \textsf{mwm}(G)
\end{align*}
Due to complementary slackness (\Cref{prop:invariant}\ref{item:slackness}), we have,
\begin{align*}
\sum_{u\in V} y(u)+\sum_{e\in E} x(e)\cdot z(e)\geqslant (1-\eps)\cdot \textsf{mwm}(G)
\end{align*}
Due to the fact that free vertex duals, at the end of the algorithm are zero (\Cref{prop:invariant}\ref{item:freeduals}), we have,
\begin{align*}
\sum_{u\in V} \sum_{e\ni u} x(e)\cdot y(u)+\sum_{e\in E} z(e)\cdot x(e)\geqslant (1-\eps)\cdot \textsf{mwm}(G)
\end{align*}
Simplifying the above expression, we have,
\begin{align*}
\sum_{e\in E} yz(e)\cdot x(e)\geqslant (1-\eps)\cdot \textsf{mwm}(G)
\end{align*}
Using \Cref{claim:forward}, we have, 
\begin{align*}
\sum_{e\in E} (1+4\eps)\cdot w(e)\cdot x(e)\geqslant (1-\eps)\cdot \textsf{mwm}(G)
\end{align*}
This shows the claim. 
\end{proof}
We now want to state the precise algorithm.
\begin{algorithm}[H] 
	\algorithmicrequire{ A bipartite multigraph $G$ with capacity and weight edge functions $\kappa$ and $w$.}\\
	\algorithmicensure{ A fractional matching $\vec{x}$ with $\sum_{e\in E}w(e)\cdot x(e)\geqslant (1-\eps)\cdot \textsf{mwm}(G,\kappa)$}
	\caption{\textsc{Weighted-Frac-Matching}($G,\kappa,\eps$)}
	\begin{algorithmic}[1]
	\Procedure{Initialization:}{}
		\State $x(e)\leftarrow 0$ for all $e\in E$ \Comment{Initializing $\vec{x}$}
		\State $y(u)\leftarrow W-\eps\cdot W$ for all $u\in L$ \Comment{Initializing left vertex duals.}
		\State $y(u)\leftarrow 0$ for all $u\in R$ \Comment{Initializing right vertex duals.}
		\State $z(e)\leftarrow 0$ for all $e\in E$ \Comment{Initializing edge duals.}
	\EndProcedure
	\State Execute scales $i=1,\cdots, \log W$ and return fractional matching $\vec{x}$.
	\State \textbf{Scale $i$}:
		\While{$y$-values of free left vertices are greater than $\frac{W-\eps\cdot W}{2^i}$}
		\For{$e_b\in G^t_x$ with $z(e)>0$}
		\State $z(e)\leftarrow z(e)-\min\set{z(e), yz(e)-w_i(e)+\delta_i}$. \Comment{Either $z(e)=0$, or $e\notin G^t_x$ as $yz(e)=w_{i}(e)-\delta_i$} \label{item:zadj} 
		\EndFor
		\State Update $G^t_x$
		\Procedure{Augmentation:}{} 
		\State Find the maximal set $\mathcal{P}$ of augmenting paths in $G^t_x$. \label{item:augpath}
		\State Augment along $\mathcal{P}$.  \label{item:augmentation}
		\State Update $\vec{x}$ and $G^t_x$.  \label{item:augmentation2}
		\EndProcedure
	\Procedure{Dual Adjustment:}{}
		\State Let $Z$ be the set of vertices reachable from free left vertices in $G^t_x$.
		\State For all ineligible edges $e_b$ directed from $R\setminus Z$ to $L\cap Z$, $z(e)\leftarrow z(e)+\delta_i$.
		\State For all $u\in L\cap Z$, $y(u)\leftarrow y(u)-\delta_i$.
		\State For all $u\in R\cap Z$, $y(u)\leftarrow y(u)+\delta_i$.
	\EndProcedure
		\EndWhile
	\Procedure{Prepare for next scale:}{}
		\State $\delta_{i+1}\leftarrow \frac{\delta_{i}}{2}$
		\State $y(u)\leftarrow y(u)+2\delta_{i+1}$ for all $u\in L$
	\EndProcedure
		\end{algorithmic}
	 \label{alg:weightedfracmatching}
\end{algorithm}

We now show that the \Cref{alg:weightedfracmatching} preserves \Cref{prop:invariant} throughout. Before this, we start with the following claim.


\begin{observation2}\label{obs:noaugmentingpaths}
After \Cref{item:augmentation} and \Cref{item:augmentation2}, the graph $G^t_x$ doesn't contain any augmenting paths. 
\end{observation2}
\begin{proof}
Suppose $G^t_x$ is the graph before the augmentation, and let $G^t_y$ be the graph after the augmentation. We want to argue that the directed edges of $G^t_y$ are a subset of the directed edges of $G^t_x$. In order to see this, consider an edge $e$, due to the eligibility criterion, either $e_f\in G^t_x$ or $e_b\in G^t_x$ but not both. \\ \\
Without loss of generality, suppose $e_f$ is present in $G^t_x$, and suppose it is augmented along and is saturated by $\vec{y}$. Then, since $G^t_y$ is a residual graph of eligible edges, and the duals of $G^t_x$ and $G^t_y$ are the same, this implies that $G^t_y$ contains neither $e_b$ nor $e_f$. On the other hand, suppose $e_f$ is not saturated by $\vec{y}$, then it appears in $G^t_y$ with capacity at most $\kappa(e)-y(e)\leqslant \kappa(e)-x(e)$, and by the same argument $e_b$ does not appear in $G^t_{y}$. A similar argument applies for $e_b$. Thus, we can conclude that the directed edges of $G^t_y$ are a subset of the edges of $G^t_x$, and therefore, any augmenting path present in the former is present in the latter as well. This contradicts the fact that we found a maximal set of augmenting paths.
\end{proof}

We now show that \textbf{granularity} property holds. 
\begin{claim2}\label{claim:domination}
Throughout \Cref{alg:weightedfracmatching}, \Cref{prop:invariant}\ref{item:gran} is always maintained. 
\end{claim2}
\begin{proof}
We will argue by induction, that for every scale $i=\set{1,\cdots, \log W}$, it holds throughout the scale. First note that it holds at the beginning of each scale, since $\delta_0$ and $W$ are assumed to be powers of $2$. Moreover, $y(u)=W-\eps\cdot W$ for all $u\in L$, and $y(v)=0$ for all $v\in R$, and $z(e)=0$ for all $e\in E$. Consider step \Cref{line:zadj}, suppose for $e=(u,v)$, $z(e)$, $y(u)$ and $y(v)$ were a multiple of $\delta_0$ before this step, then after this step, $z(e)$ is still a multiple of $\delta_0$. This is because $w_i(e)$ is a multiple of $\delta_0$. Similarly, in the dual adjustment procedure, we increment and decrement the duals by $\delta_0$, thus, if they were multiples of $\delta_0$, then they would be multiples of $\delta_0$ even after this step. Finally, when we prepare for the next scale, we make sure that the duals are a multiple of $\delta_1$. We can similarly make such an argument for all $i\geqslant 2$. 
\end{proof}
We now show \textbf{near domination}. 
\begin{claim2}
\Cref{prop:invariant}\ref{item:domin} is maintained throughout the run of \Cref{alg:weightedfracmatching}.
\end{claim2}
\begin{proof}
Observe that \Cref{prop:invariant}\ref{item:domin} holds at the beginning of the $0$th scale. Assume it holds at the beginning of scale $i$. We will show that it holds at the beginning of scale $i+1$ as well. We do so by showing that during scale $i$, each time we change duals, we preserve \textbf{domination}. Consider \Cref{item:zadj}, observe that we modify $z(e)$ in this case if $e_b$ is an eligible edge. This implies that $yz(e_b)\geqslant w(e_b)$. Consequently, after the adjustment to $z(e)$, \Cref{prop:invariant}\ref{item:domin} still holds since after the dual adjusment, $yz(e_b)\geqslant w(e_b)-\delta_i$. \\ \\
Consider now the dual adjustment step, observe that for all edges $e\in E(L\cap Z,R\cap Z)$ and $e\in E(L\setminus Z, R\setminus Z)$ remain eligible after the dual adjustment step since their $yz(e)$ value remains the same. Now, observe there are no eligible edges $e_f\in E(L\cap Z, R\setminus Z)$. This implies that we have $yz(e_f)>  w_i(e_f)-\delta_i$. However, since vertex and edge duals were all multiples of $\delta_i$ (see \Cref{claim:domination}), we can conclude that $yz(e_f)\geqslant w_i(e_f)$. Thus, after dual adjustment step, we have $yz(e_f)\geqslant w_i(e_f)-\delta_i$. Thus, \Cref{prop:invariant}\ref{item:domin} holds after this step. Consider $e_b\in E(R\setminus Z, L\cap Z)$, if this edge is eligible then, after dual adjustment, $yz(e_b)-w_i(e_b)\geqslant -\delta_i$, and domination property still holds. On the other hand, if $e_b$ is ineligible then, we also increment $z(e_b)$, and consequently, $yz(e)$ remains the same, and we still have domination property. \\ \\
Next, consider edges between $R\cap Z$ and $L\setminus Z$, for such edges, since we only increase the duals of $R\cap Z$, and duals of $L\setminus Z$ are unaffected, we still have the domination property. \\ \\
Consider the time when we are preparing for the next scale. In this case, for all vertices $u\in L$, we increment $y(u)$ by $2\delta_{i+1}$. Thus, we have,
\begin{align*}
yz(e)\geqslant w_{i}(e)-\delta_i+2\delta_{i+1}\geqslant w_{i}(e) \geqslant w_{i+1}(e)-\delta_{i+1}
\end{align*}
This proves the claim
\end{proof}
\begin{observation2}
Note that \Cref{prop:invariant}\ref{item:freeduals} holds, since due to property \Cref{obs:noaugmentingpaths}, all free vertices in $L$ are a subset of $L\cap Z$ and all free vertices of $R$ are contained in $R\setminus Z$. Thus, at the beginning of the algorithm, free vertex duals in $R$ are $0$, and they remain $0$ during the rest of the algorithm, since duals of vertices in $R\setminus Z$ are not modified.
\end{observation2}

We now show the complementary slackness property.

\begin{claim2}
Note that \Cref{prop:invariant}\ref{item:slackness} holds throughout the algorithm.
\end{claim2}
\begin{proof}
To see this, observe that $z(e)$ is only incremented when $e_b\in E(R\setminus Z, L\cap Z)$ and $e_b$ is ineligible. Note since $e_b$ is ineligible, it must be the case that $yz(e)-w_i(e)=-\delta_i$. Thus, $e_f$ if it exists, must be eligible, however this would contradict the fact that there can't be eligible forward edges between $R\setminus Z$ and $L\cap Z$. Consequently, $e_f$ doesn't exist, and therefore, we have $x(e)=\kappa(e)$. On the other hand, observe that if $e_b$ with $z(e)>0$ is added to the eligible graph, then it is assured that $z(e)=0$ before $x(e)$ is reduced from $\kappa(e)$ to something lower. This concludes the proof. 
\end{proof}

We now show the tightness property. 

\begin{claim2}
Note that \Cref{prop:invariant}\ref{item:tightness} holds throughout the algorithm. 
\end{claim2}
\begin{proof}
Observe that this holds during the augmentation phase. We now want to show this holds during the dual adjustment phase as well. Note that, as before the edges in $E(L\cap Z, R\cap Z)$ and $E(L\setminus Z, R\setminus Z)$ are unaffected. On the other hand, backward edges in $E(L\cap Z, R\setminus Z)$ will also not have their tighness property affected since $L\cap Z$ is only being decremented. On the other hand, there are no eligible backward edges in $E(R\cap Z, L\setminus Z)$, so the only backward edges are ineligible, which means that $yz(e)=w_{i}(e)-\delta_i$. Thus, after dual adjustment, for such edges, $yz(e)=w_i(e)$, and therefore tightness constraint is satisfied. Finally consider the dual adjustment step that prepares for the next scale, for any backward edge of type $j$, after this step, we have,
\begin{align*}
yz(e)\leqslant w_{i}(e)+2\cdot (\delta_j-\delta_i)+2\delta_{i+1}\leqslant w_{i+1}(e)+2\cdot (\delta_j-\delta_{i+1})
\end{align*}
\end{proof}

Finally, we show the following property of the eligible graph, which will enable us to bound the runtime of the algorithm.

\begin{lemma2}

\end{lemma2}

\subsection{Runtime}

\begin{observation2}
Observe that scale $i$ starts with free duals being $\frac{W-\eps\cdot W}{2^{i+1}}$ and ends with free duals being $\frac{W-\eps\cdot W}{2^{i+1}}$, and each iteration in scale $i$ ends with reducing the dual variables with $\delta_i$, so the total number of iterations in each scale is at most $\frac{1}{\eps}$
\end{observation2}

\begin{observation2}
\Cref{item:zadj} takes time $O(m)$ and is executed once per iteration in the $i$th scale. 
\end{observation2}

\begin{observation2}
\Cref{item:augpath} takes $O(m\cdot \frac{\log n}{\eps})$ time and is executed once per iteration in the $i$th scale.
\end{observation2}

\begin{observation2}
The dual adjustment steps take $O(m)$ time per iteration per scale. 
\end{observation2}

\begin{claim2}
The total runtime of the algorithm is $O(m\cdot \frac{\log n}{\eps}\cdot \log W)$. 
\end{claim2}